% Job posting: https://ca.linkedin.com/jobs/view/embedded-artificial-intelligence-architect-at-mda-space-4414774301
% =====================================================================
%  CV - Nishal Stanislaus Silva, PhD
%  Target: MDA Space - Embedded Artificial Intelligence Architect (Technical Lead)
%          Satellite Systems team, Sainte-Anne-de-Bellevue, QC
%  Variant: V3 (Edge / Embedded). Embedded ML deployment under device
%           budgets + simulator-driven training data + benchmark
%           methodology + time-series prediction framing. Honest on the
%           architect-level and satellite-comms stretch (see changelog).
% =====================================================================

\documentclass[11pt]{article}
\usepackage[left=0.7in,top=0.7in,right=0.7in,bottom=0.7in]{geometry}
\usepackage{enumitem}
\usepackage[hidelinks]{hyperref}
\usepackage{titlesec}
\usepackage{array}
\usepackage{setspace}
\usepackage{needspace}
\usepackage[T1]{fontenc}
\usepackage{newtxtext,newtxmath}
\ifdefined\pdfgentounicode
  \input{glyphtounicode}
  \pdfgentounicode=1
\fi
\setstretch{1.05}
\pagenumbering{gobble}
\titleformat{\section}{\Large\scshape}{}{4em}{}[\vspace{-0.7em}\rule{\linewidth}{0.4pt}\vspace{-0.4em}]
\titlespacing*{\section}{0pt}{1em}{0.4em}
\newenvironment{cvrole}[5]{%
  \par\noindent\textbf{\large #1} | \textit{\small #2, #3}\hfill \textit{\small #4 - #5}\par
  \begin{itemize}[leftmargin=2em, itemsep=-0.3em, topsep=-0.5em]%
}{%
  \end{itemize}\par\noindent\mbox{}\par\vspace{0.1em}%
}
\newcommand{\cvName}{Nishal Stanislaus Silva}
\newcommand{\cvEmail}{nishal.silva@hotmail.com}
\newcommand{\cvPhone}{+1 613 200 2030}
\newcommand{\cvCity}{Toronto, ON}
\newcommand{\cvSite}{https://nishal.xyz}
\newcommand{\cvLinkedIn}{https://www.linkedin.com/in/nishal-silva}
\newcommand{\cvGitHub}{https://github.com/ns2max}
\newcommand{\cvStatus}{Canadian Permanent Resident, Eligible to work in Canada without sponsorship}
\newcommand{\cvTagline}{ML Research Engineer | Real-Time \& Edge Inference | Audio, Sensor \& Time-Series}

\begin{document}
\pagestyle{empty}
\begin{center}
    {\LARGE \scshape \textbf{\cvName}}{\large \textit{, PhD}}\\[1pt]
    {\small \cvTagline}\\[1pt]
    {\footnotesize\textit{\cvStatus}}\\[1pt]
    {\small
    \href{mailto:\cvEmail}{\cvEmail} \,\,|\,\, \cvPhone \,\,|\,\, \cvCity
    \,\,|\,\, \href{\cvSite}{nishal.xyz} \,\,|\,\, \href{\cvLinkedIn}{linkedin.com/in/nishal-silva}
    \,\,|\,\, \href{\cvGitHub}{github.com/ns2max}}
\end{center}

\section*{Professional Summary}

Machine learning research engineer who deploys and prototypes ML models on embedded targets under hard compute, power, memory and latency constraints, with a PhD (University of Trento, 2025) and 10+ years across industry and academia. I select model architectures to fit an embedded budget, define metrics and evaluate model performance, robustness and resource utilization on the target platform, and design simulators and rule-based generators that produce representative training data when real data is scarce. My published detector runs inference in 14 ms on a Raspberry Pi 4 at 31.4\% CPU, F1 0.76, 74$\times$ faster than the 1038 ms DTW baseline it replaced (JAES 2026). Time-series and sequence modelling (RNN, LSTM, GRU) is my core: forecasting and pattern detection over streaming signals, the same shape as network-traffic prediction. I own streaming pipelines end to end (acquisition, DSP feature front-ends, training and evaluation under domain shift, on-device deployment on Raspberry Pi and ARM embedded Linux with monitoring), have curated four open datasets with frozen evaluation protocols (9,800+ recordings, 70+ participants), published 10+ peer-reviewed papers in JAES, IEEE and ACM venues, and mentored engineers and graduate students.

\section*{Core Competencies}

\textbf{ML Engineering:} deep learning for audio and sensor streams, time-series and sequence modelling (RNN, LSTM, GRU) for forecasting and pattern detection, CNNs, hierarchical state machines, embedding-based similarity matching, anomaly detection, offline training and online adaptation, generative models (VAE, diffusion) for scarce-label regimes, TensorFlow/Keras, PyTorch, Python, production C/C++17\\
\textbf{Edge \& System Optimization:} deploying and prototyping ML on embedded targets (Raspberry Pi, ARM) under compute, power, memory and latency budgets, selecting model architectures to fit those budgets, defining metrics and evaluating model performance, robustness and resource utilization on target platforms, designing simulators and rule-based data generators for representative training data, dataset curation with frozen train/validation/test protocols, model quantization for deployment, real-time streaming inference under sub-30 ms budgets, embedded Linux, C++17 real-time engines, latency-aware edge-to-server partitioning\\
\textbf{Research Engineering:} benchmark and dataset design, frozen-protocol evaluation, ablation studies, trade-off analysis of accuracy against compute and memory, technical documentation of architecture, design, verification and deployment, written experiment records including negative results, 10+ papers, peer review and programme-committee service, mentoring, communication with non-ML domain and hardware experts

\section*{Technical Stack}

\textbf{Languages:} Python, C++17, C, MATLAB, JavaScript/TypeScript, Shell\\
\textbf{ML \& AI:} TensorFlow, Keras (embedded deployment), PyTorch, scikit-learn, librosa, torchaudio; RNN/LSTM/GRU, CNN, embedding matching, VAE and diffusion prototypes; model compression and quantization for deployment; NumPy, Pandas, SciPy\\
\textbf{Edge \& Systems:} Raspberry Pi (production), ARM CPU latency and memory profiling, embedded Linux, real device traces, simulators and rule-based data generators, rule-based and generative synthetic training data, Elk Audio OS, JUCE, VST, OSC, FFTW, libsndfile, IMU/I\textsuperscript{2}C, PLC interfacing, OpenCV; familiarity with ONNX, TFLite\\
\textbf{Data \& Dev:} AWS (EC2, S3, ECS, MWAA, RDS), SQL, ETL pipelines, Docker, Airflow, Jenkins, Git, Linux/UNIX, pytest, CI/CD, REST APIs, Jira, Confluence

\needspace{5\baselineskip}
\section*{Portfolio \footnotesize {\normalfont{(full portfolio: \href{https://nishal.xyz/research}{\textit{nishal.xyz/research}})}}}
\begin{itemize}[leftmargin=1.2em, itemsep=-0.25em]
    \item \textbf{Nebula} \href{https://github.com/ns2max/nebula}{\textit{(github.com/ns2max/nebula)}}: C++17 audio-feature library with per-feature ARM / Raspberry Pi latency benchmarks across time-domain, spectral, time-frequency and cepstral features.
    \item \textbf{Hot Licks} \href{https://youtu.be/gkOqfOT8zXM}{\textit{(youtu.be/gkOqfOT8zXM)}}: real-time audio pattern detection on Raspberry Pi; MIDI Innovation Awards 2023 finalist.
\end{itemize}

\section*{Education}
\textbf{Ph.D - Information Engineering and Computer Science} \textit{\small{ - University of Trento, Italy}} \hfill \textit{Jan 2025}\\[2pt]
\noindent\textit{\small Thesis: Embedded Real-time Musical Pattern Detection for Smart Musical Instruments}\\[3pt]
\noindent\textbf{M.Sc - Telecommunication and Electronic Engineering} \textit{\small{ - Sheffield Hallam University, UK}} \hfill \textit{Aug 2018}\\[2pt]
\noindent\textit{\small Thesis: On Musical Onset Detection via the S-Transform (Asilomar 2018)}\\[3pt]
\noindent\textbf{B.Eng (Hons) - Electronic Engineering} \textit{\small{ - Sheffield Hallam University, UK}} \hfill \textit{Mar 2014}

\section*{Research and Work Experience}

\begin{cvrole}{Independent Research \& Open-Source}{Self-directed}{Toronto, Canada}{Jan 2026}{Present}
    \item Nebula: C++17 audio-feature-extraction library with per-feature ARM / Raspberry Pi latency benchmarks across time-domain, spectral, time-frequency and cepstral features.
    \item MIRaaS paper accepted at IEEE IS\textsuperscript{2} 2026; Smart Drums paper in press at JAES; side projects built with AI-assisted workflows including Claude Code.
\end{cvrole}

\begin{cvrole}{Postdoctoral Researcher}{University of Trento}{Trento, Italy}{Jan 2025}{Dec 2025}
    \item Owned the end-to-end streaming-audio and IMU/gesture ML pipeline: acquisition, DSP feature front-end, training and evaluation, and low-latency on-device deployment (Raspberry Pi 4, embedded Linux, Elk Audio OS, VST, OSC) with monitoring; EU Horizon EIC Pathfinder project MUSMET.
    \item Delivered real-time inference at 14 ms on a Raspberry Pi 4, F1 0.76, 74$\times$ faster than the 1038 ms DTW baseline, at 31.4\% CPU within an embedded thermal and memory budget (JAES 2026).
    \item Selected model architectures for the device budget and built benchmark suites with frozen protocols: baselines and ablations (RNN vs DTW, audio vs MIDI, pattern-set sizes 1/3/10) quantifying accuracy, latency, CPU and memory trade-offs on the target platform, with a written record of experiments including negative results.
    \item Built rule-based synthetic-variation generators (roughly 10k variants per pattern) and VAE/diffusion synthetic data for scarce-label conditions; curated frozen train/validation/test splits.
\end{cvrole}

\needspace{5\baselineskip}
\begin{cvrole}{Visiting Researcher}{McGill University (IDMIL / CIRMMT)}{Montreal, Canada}{May 2024}{Aug 2024}
    \item Fused inertial (IMU) and audio streams through a hierarchical finite state machine on an instrument-mounted Raspberry Pi 4, reaching F1 0.78 across a 500-event corpus (IEEE IS\textsuperscript{2} 2025).
    \item Profiled compute and power for real-time feasibility on self-powered embedded hardware.
\end{cvrole}

\needspace{5\baselineskip}
\begin{cvrole}{Technical Consultant}{Forestpin (Pvt) Ltd}{Colombo, Sri Lanka}{Aug 2020}{Feb 2021}
    \item Designed and deployed ML and statistical pipelines for anomaly detection and risk scoring on large structured enterprise datasets; SQL and Python backend scoring services in production.
\end{cvrole}

\needspace{5\baselineskip}
\begin{cvrole}{Research Engineer}{MAS Holdings (Pvt) Ltd}{Colombo, Sri Lanka}{Jan 2015}{Jul 2020}
    \item Built end-to-end computer-vision and sensor systems for manufacturing at scale, integrating camera, sensor and IoT streams into production workflows; operational efficiency improved by up to 300\%.
    \item Shipped automated multilingual care-label QC (OpenCV template and feature matching, conveyor with PLC reject); inspection time cut 99.5\%, with an explainable defect overlay for line operators.
    \item Built real-time ingestion and ETL for high-frequency IoT and operational time-series across sites in South Asia, North America and Africa; C++/Python inference engines on embedded hardware; introduced CI/CD and code review; mentored engineers and ran group-wide training.
\end{cvrole}

\section*{Datasets \footnotesize {\normalfont{(open access, Zenodo)}}}
Four open musical-pattern datasets (DoMP, DoPP, DoDP, DoDP2): 9,800+ recordings, 70+ musicians, with a unified artist-ID namespace so cross-performer and cross-device splits stay honest for domain-shift and on-device evaluation.

\section*{Selected Publications \footnotesize {\normalfont{(full list: \href{https://nishal.xyz/publications}{\textit{nishal.xyz/publications}})}}}
\begin{itemize}[leftmargin=1.2em, itemsep=-0.25em]
    \item Silva, N. and Turchet, L. \textit{Real-Time Audio Pattern Detection for Smart Musical Instruments}. Journal of the Audio Engineering Society, Vol.~74, No.~3, Mar.~2026. \textit{(F1 0.76 at 14 ms on Raspberry Pi 4, 31.4\% CPU, 74$\times$ faster than baseline)}
    \item Silva, N. and Turchet, L. \textit{Smart Drums: Comparing Audio and MIDI in Embedded Real-Time Drum Pattern Recognition}. Journal of the Audio Engineering Society, 2026 (in press).
    \item Silva, N. and Turchet, L. \textit{Music Information Retrieval as a Service: Latency Characterization of an Edge-to-Server Architecture for Near-Real-Time Audio Analysis}. IEEE International Symposium on the Internet of Sounds, 2026 (accepted).
    \item Silva, N. and Turchet, L. \textit{A Structural Similarity Index Based Method to Detect Symbolic Monophonic Patterns in Real-Time}. 25th International Conference on Digital Audio Effects (DAFx), Sep.~2022.
    \item Silva, N., Weeraddana, P.C. and Fischione, C. \textit{On Musical Onset Detection via the S-Transform}. 52nd Asilomar Conference on Signals, Systems, and Computers, Oct.~2018.
\end{itemize}

\section*{Selected Projects \footnotesize {\normalfont{(full portfolio: \href{https://nishal.xyz/research}{\textit{nishal.xyz/research}})}}}
\begin{itemize}[leftmargin=1.2em, itemsep=-0.25em]
    \item \textbf{Real-time streaming pattern detection under a device budget:} 30 ms windows / 10 ms hop, MFCC front-end into a stacked RNN, trained on rule-based synthetic variations; Raspberry Pi 4 at 14 ms, 31.4\% CPU, F1 0.76, against a DTW baseline at 1038 ms.
    \item \textbf{Simulator-driven training data:} rule-based generators and VAE/diffusion synthesis producing roughly 10k representative variants per pattern for low-label regimes, with frozen evaluation protocols across human, synthetic and public reference sets.
    \item \textbf{Multimodal sensor fusion on embedded hardware:} IMU and audio fusion via a hierarchical state machine on a self-powered instrument-mounted Raspberry Pi 4; F1 0.78 on a 500-event corpus, with compute and power profiling.
\end{itemize}

\vspace{-0.5em}
\enlargethispage{5\baselineskip}
\section*{Highlights \& Recognition}
\begin{itemize}[leftmargin=1.2em, itemsep=-0.25em]
    \item \textbf{Evaluation methodology:} designed frozen-protocol benchmark suites with cross-condition evaluation tiers and a written record including negative results, quantifying accuracy, latency, CPU and memory trade-offs on the deployment platform.
    \item \textbf{MIDI Innovation Awards} Finalist (Sep 2023); reviewer, IEEE Access and the J.\ National Science Foundation of Sri Lanka; programme committee, IEEE IS\textsuperscript{2} and IEEE I3DA.
    \item \textbf{Languages:} English (native); French (basic, improving).
\end{itemize}

\end{document}
